\subsection{马施克定理}

定理 1. —— 设群 G 是有限的。设 M 是一个 K{[}G{]}-模，N 是 M 的一个 K{[}G{]}-子模。假设 N 是 K-模 M 的一个直因子，且 \textbar G\textbar{} 在 K 中可逆。那么 N 是 K{[}G{]}-模 M 的一个直因子。

设 p 是 M 中以 N 为像的一个 K-线性投影算子。我们如下定义 K-模 M 的一个自同态 q：

\[
q (x) = | \mathrm{G} | ^ {- 1} \sum_ {g \in \mathrm{G}} g p (g ^ {- 1} x)
\]

对每个 \(x \in M\) 成立。由于 N 在 G 的作用下稳定，且 p 在 N 上诱导恒同，我们看到 q 把 M 映入 N 且在 N 上诱导恒同。

于是 K-模 M 是 q 的像 N 与 q 的核的直和。\(g q(x) = q(gx)\)（对每个 \(x \in M\) 与 \(g \in G\)），故 q 的核是 M 的一个 K{[}G{]}-子模。这就证明了 N 是 K{[}G{]}-模 M 的一个直因子。

推论 1（马施克）. —— 设群 G 有限且 K 是一个交换域。代数 K{[}G{]} 是半单的，当且仅当域 K 的元素 \textbar G\textbar{} 不为零。

设 \(|\mathrm{G}| \neq 0\)。由定理 1，每个 K{[}G{]}-子模（\(\mathrm{K}[\mathrm{G}]_s\) 的子模）都是直因子。因此 K{[}G{]} 是半单的。

反之，设 \(|G|\) 为零，并用 \(\varepsilon\) 表示元素 \(\sum_{g\in G} g\)（K{[}G{]} 的中心的元素）。我们有 \(\varepsilon \neq 0\) 但 \(\varepsilon^{2} = |G|\varepsilon = 0\)，故 K{[}G{]} 不是半单的（VIII, 第 153 页，命题 3, a) 与 VIII, 第 157 页，注 1）。

推论 2. —— 设群 G 有限且 \(|G|\) 在 K 中可逆。K{[}G{]}-模的一个正合列分裂，当且仅当它作为 K-模的正合列分裂。

给定 K{[}G{]}-模的一个正合列

\[
0 \longrightarrow \mathrm{M} ^ {\prime} \stackrel {f} {\longrightarrow} \mathrm{M} \longrightarrow \mathrm{M} ^ {\prime \prime} \longrightarrow 0,
\]

只需对态射 f 的像应用定理 1。

推论 3. —— 设群 G 有限且 \(|G|\) 在 K 中可逆。一个 K{[}G{]}-模是投射的，当且仅当它作为 K-模是投射的。

设 P 是一个 K{[}G{]}-模。若 P 是自由 K{[}G{]}-模 M 的一个直因子，则它更是自由 K-模 M 的一个直因子。其逆由推论 2 得出，参见 II, §2, No.~2, 第 231 页，命题 4, d)。

推论 4. —— a) 设群 G 有限且 K 是一个特征为 0 的交换域。G 的两个有限维线性表示同构，当且仅当它们有相同的特征标。

\begin{enumerate}
\def\labelenumi{\alph{enumi})}
\setcounter{enumi}{1}
\tightlist
\item
  设 K 是一个特征为素数 p 的完满域，且群 G 有限、基数与 p 互素。G 的一个有限维线性表示的特征标为零，当且仅当该表示同构于 p 个两两同构的表示的直和。
\end{enumerate}

在该推论的假设下，G 的每个有限维线性表示都是半单的。于是该推论由 VIII, 第 384 页的推论得出。

推论 5. —— 设群 G 有限且 K 是一个满足 \(|G| \neq 0\) 的交换域。设 \(\pi\) 与 \(\pi'\) 是 G 的有限维线性表示。那么 \(\pi\) 与 \(\pi'\) 同构，当且仅当对 G 中的每个 g，自同态 \(\pi(g)\) 与 \(\pi'(g)\) 有相同的特征多项式。

这由 VIII, 第 386 页的推论 1 得出。

